\documentclass[11pt]{article}

\usepackage[margin=1.1in]{geometry}
\usepackage{amsmath}
\usepackage{amssymb}
\usepackage{amsthm}
\usepackage{booktabs}
\usepackage{array}
\usepackage[T1]{fontenc}
\usepackage[hidelinks]{hyperref}

\theoremstyle{plain}
\newtheorem{theorem}{Theorem}[section]
\newtheorem{lemma}[theorem]{Lemma}
\newtheorem{proposition}[theorem]{Proposition}
\newtheorem{corollary}[theorem]{Corollary}
\theoremstyle{definition}
\newtheorem{definition}[theorem]{Definition}
\theoremstyle{remark}
\newtheorem*{remark}{Remark}

\numberwithin{equation}{section}

\newcommand{\N}{\mathbb{N}}
\newcommand{\Z}{\mathbb{Z}}
\newcommand{\Q}{\mathbb{Q}}
\newcommand{\R}{\mathbb{R}}
\DeclareMathOperator{\Ac}{A}

\title{Machine-Verified Resolutions of Four OEIS Conjectures}

\author{Haoyu Chen\thanks{%
Disclosure of the production pipeline. Every mathematical step reported in this paper was
produced by automated systems; no human supplied a definition, a lemma, a proof idea, or a
correction. The proof of Theorem~\ref{thm:a108211} was found by Claude (Fable 5) after an
earlier attempt was refuted in adversarial review. Theorem~\ref{thm:a114831} was solved twice
independently: once by Claude (Fable 5) and once, blind to the first solution, by OpenAI
GPT-5.6 through the \texttt{codex} command-line interface; the second derivation is the one
reproduced here. Theorem~\ref{thm:a100434} was found by Claude (Fable 5) after a numerical
reconnaissance run refuted the statement as literally transcribed in an existing formal
corpus. Theorem~\ref{thm:a114362} was solved by Claude (Fable 5), refuted in adversarial review
by two independent reviewers who located the same defect, and then repaired by GPT-5.6
(``Pro'' effort) --- the reviewer that had supplied the counterexample --- with the repair itself
independently re-validated by two further reviewers before being accepted
(Section~\ref{sec:a114362}). Adversarial review was carried out in separate sessions by GPT-5.6
(``Pro'' effort, web interface), Qwen3.8-Max (web interface), and GPT-5.6 through \texttt{codex},
under a protocol requiring the reviewing model to differ from the solving model, except for the
one repair step just noted. All four theorems were then formalized in Lean~4 against
\texttt{mathlib}; the Lean compiler is the final arbiter of every claim below, and the files are
listed in Section~\ref{sec:formal}.
Orchestration --- problem selection, task routing, review scheduling, and the drafting of
this paper --- was likewise automated. Responsibility for what follows rests with the reader
who checks the Lean files.}}

\date{16 August 2026}

\begin{document}

\maketitle

\begin{abstract}
We resolve four conjectures recorded in the On-Line Encyclopedia of Integer Sequences and
formalize each resolution in Lean~4. First, we prove a conjecture of C.~Kimberling (2014,
A108211): for every integer $n \ge 1$,
$\bigl\lfloor \bigl(\tfrac{1}{4n} - \log 2 + \sum_{k=n+1}^{2n} \tfrac1k \bigr)^{-1}\bigr\rfloor
= 16n^2+1$. The proof is entirely certificate-based: a rational telescoping majorant locates
the tail of the paired alternating harmonic series to within $O(n^{-7})$, comfortably inside a
target window of width $\Theta(n^{-4})$, and all five auxiliary inequalities are reduced to
polynomials whose coefficients, after the substitution $x = 1+t$, are nonnegative; the
coefficient lists are reproduced in full in Appendix~\ref{app:certificates}. Second, we answer
the asymptotic question posed with A114831 (J.~V.~Post, 2006): the sequence defined by
$a(1)=1$, $a(2)=2$, $a(n) = a(n-1) + \lfloor 2a(n-1)a(n-2)/(a(n-1)+a(n-2))\rfloor$ satisfies
$a(n+1)/a(n) \to \sqrt3$, with the explicit rate
$|a(n)/a(n-1) - \sqrt3| \le (2-\sqrt3)2^{-(n-2)} + 4/n + 2^{-(n-2)/2}$ for $n \ge 4$. An
independent, essentially simultaneous kernel-checked Lean proof of this limit was posted by
K.~Kitamura shortly before ours; we make no claim to priority for this theorem (see the Prior
work remark in Section~\ref{sec:a114831}). Third, we
prove the three linear-recurrence identities conjectured by C.~Dement (2004, A100434), in the
sign-corrected form that the source's own ``apart from signs'' convention dictates; the literal
transcription is refutable at $n=0$. Fourth, we prove the asymptotic conjectured by
T.~Ordowski (2022, A114362 and A348829): writing $t(n) := \zeta(2n)/\zeta(n)^2$, the quantity
$y(n) := (1-t(n))/(1+t(n))$ satisfies $y(n) = 2^{-n}+3^{-n}+5^{-n}+7^{-n}+O(11^{-n})$; we obtain
the explicit bound $\bigl|y(n)-(2^{-n}+3^{-n}+5^{-n}+7^{-n})\bigr| \le 24\cdot11^{-n}$ for $n\ge2$
from the Euler product for $\zeta$ by an exact identity for a nested composition of M\"obius maps
together with an elementary tail estimate; the Lean development proves the same $O(11^{-n})$
statement with an explicit, weaker witness constant of $56$. This theorem's first proof was
independently refuted by two reviewers before a second, corrected proof was found and
dual-reviewed as valid. Each theorem is accompanied by a Lean~4 development whose
axiom audit reports dependence only on \texttt{propext}, \texttt{Classical.choice}, and
\texttt{Quot.sound}. As a by-product we report two mis-formalizations in the
\texttt{formal-conjectures} corpus (A109074 and A100434).
\end{abstract}

\section{Introduction}

\subsection{Four conjectures}

The On-Line Encyclopedia of Integer Sequences \cite{oeis} carries, alongside its sequence data,
a large stock of conjectural comments contributed over three decades. Most are elementary in
the sense that no deep machinery is expected to be needed; many are nevertheless unproved,
because nobody has taken the time. This paper resolves four of them.

\begin{enumerate}
\item[(A108211)] Clark Kimberling observed on 9 September 2014 \cite{kimberling} that the sequence
$a(n) = 16n^2+1$ appears to be given by
\[
  a(n) \;=\; \left\lfloor \frac{1}{\frac{1}{4n} - \log 2 + \frac{1}{n+1} + \frac{1}{n+2}
    + \cdots + \frac{1}{2n}} \right\rfloor .
\]
This is Theorem~\ref{thm:a108211} below. It is the substantive result of the paper: the
quantity being inverted is a difference of two quantities that agree to relative order $n^{-1}$,
and the floor is pinned only if the difference is located inside an interval of width
$\Theta(n^{-4})$.

\item[(A114831)] Jonathan Vos Post contributed on 19 February 2006 the sequence
\[
  1,\; 2,\; 3,\; 5,\; 8,\; 14,\; 24,\; 41,\; 71,\; 122,\; 211,\; \dots
\]
given by $a(1)=1$, $a(2)=2$ and
\[
  a(n) \;=\; a(n-1) + \left\lfloor \frac{2\,a(n-1)\,a(n-2)}{a(n-1)+a(n-2)} \right\rfloor
  \qquad (n \ge 3),
\]
together with the
question: \emph{what is this sequence, asymptotically?} The heuristic answer $\sqrt3$ is
immediate --- if the ratio of consecutive terms converges to $L$ then $L = 1 + 2/(L+1)$, so
$L^2 = 3$ --- but the existence of the limit is the whole content of the question.
Theorem~\ref{thm:a114831} settles it with an explicit error bound; an independent, essentially
simultaneous kernel-checked Lean proof of the same limit was found by another author, and we
cede priority for it (Section~\ref{sec:a114831}).

\item[(A100434)] Creighton Dement contributed on 18 December 2004 \cite{dement}, in the context of
his ``floretion'' calculus, a family of auxiliary sequences attached to the linear recurrence
$x(n+4) = -6x(n+2) - x(n)$ and conjectured three identities among them.
Theorem~\ref{thm:a100434} proves them. As discussed in Section~\ref{sec:a100434}, the identities
must be read with the sign convention that Dement's comment announces but does not carry through
its own notation; under the literal reading they fail at $n=0$.

\item[(A114362)] Thomas Ordowski commented on 13 November 2022 \cite{ordowski} that, writing
$t(n) := \zeta(2n)/\zeta(n)^2$,
\[
  \frac{1-t(n)}{1+t(n)} \;=\; \frac{1}{2^n} + \frac{1}{3^n} + \frac{1}{5^n} + \frac{1}{7^n}
    + O\!\left(\frac{1}{11^n}\right).
\]
This is Theorem~\ref{thm:a114362} below. The identity
$\zeta(2n)/\zeta(n)^2 = \prod_p (1-p^{-n})/(1+p^{-n})$ turns the asymptotic into a statement about
an infinite product of M\"obius transformations of the primes; the argument is a case study in
extracting an explicit error bound, polynomial in $1/11^n$, from a convergent Euler product, and
is the one theorem in this paper whose first proof was independently refuted by two reviewers
before being repaired.
\end{enumerate}

\subsection{Verification methodology}

All four of these statements had been transcribed into Lean~4 by the
\texttt{formal-conjectures} project \cite{formalconjectures}, a corpus of open conjectures
stated formally with \texttt{sorry} in place of a proof. That corpus supplied both our problem
list and, for A108211, A114831, and A114362, the exact formal statement to be proved, which
removes the usual risk that an autoformalization proves something other than what was asked;
A100434 is the exception, discussed in Section~\ref{sec:upstream}.

Our working protocol had three independent gates, and a claim was recorded as resolved only
after all three were passed.

\emph{Adversarial review.} A completed argument was submitted, in separate sessions and without
the solver's reasoning trace, to models from a different vendor than the solver, with
instructions to classify every objection as \emph{critical error} or \emph{justification gap}
and to return a verdict. This gate did real work. The first proof of Theorem~\ref{thm:a108211}
was killed by it: the reviewer produced an explicit counterexample at $m=2$ to a comparison
inequality on which the argument rested, and, more damagingly, showed that even had the
inequality been true its precision --- an uncertainty of order $3/(256n^4)$ against a permitted
window of order $1/(256n^4)$ --- would have been insufficient. The proof presented in
Section~\ref{sec:a108211} is the replacement, designed to have precision $O(n^{-7})$ against the
same window. Two further rounds of review, in separate sessions, returned \textsc{valid} --- one
recomputing every certificate symbolically from scratch, the other (GPT-5.6 at ``Pro'' effort)
additionally supplying, unprompted, the exact rational $12743/21656250$ for the tightest of them
at $n=1$. A third round, by Qwen3.8-Max, returned \textsc{incomplete} with eight presentation
objections, all of the form
``the certificate is asserted, not exhibited''. Section~\ref{sec:a108211} and
Appendix~\ref{app:certificates} are written to answer those eight objections: every certificate
polynomial is displayed, every quantity is an exact rational, the pairing step is justified by
an explicit subsequence argument, and the strictness of the inequality that survives the
infinite summation is obtained by retaining the first-term deficit rather than by appeal to
term-wise strictness.

A parallel failure recurred with Theorem~\ref{thm:a114362}: an initial proof asserted uniform
pointwise bounds on a nested composition of M\"obius maps that fail already at $n=2$; two
reviewers, working independently, isolated the same defect, one of them supplying an exact
rational counterexample. Section~\ref{sec:a114362} records the repair and its own independent
re-validation. Taken together, these are the two cases in this paper where the adversarial gate,
and not the solver, is the reason the final theorem is true.

\emph{Formal verification.} Each proof was then formalized in Lean~4 against
\texttt{mathlib} \cite{lean4,mathlib}, and the resulting theorem's axiom dependencies were
audited. The Lean compiler is the arbiter: no informal step in this paper is load-bearing that
is not also present in the corresponding \texttt{.lean} file.

\emph{Statement fidelity.} Because a formal proof of the wrong statement is worthless, the final
theorem statement of each Lean file was compared, by hand-transcription, against the OEIS
comment it is supposed to formalize. The comparison is recorded in Section~\ref{sec:formal};
where our statement differs from \texttt{formal-conjectures}' (as it must for A100434, and as it
does cosmetically for A114831), we say exactly how.

\subsection{Two mis-formalizations found as a by-product}
\label{sec:upstream}

Running numerical reconnaissance against formal statements before attempting to prove them
turned up two files in \texttt{formal-conjectures} whose Lean statement does not express the
OEIS conjecture it cites, and which are in fact refutable.

\begin{itemize}
\item The file \texttt{OEIS/109074.lean} defines
$b(n) = \binom{3n}{n}/(2n+1)$ and asserts
$\binom{6n-2}{2n} / \bigl(2\binom{4n-1}{2n}\bigr) = b(n+1)/b(n)$ for $n \ge 1$. Two errors
compound. First, $\binom{3n}{n}/(2n+1)$ is A001764 (ternary tree numbers
$1,1,3,12,55,\dots$), not A005156 (vertically symmetric alternating sign matrices,
$1,1,3,26,646,\dots$) as the surrounding documentation states. Second, the OEIS conjecture uses
the ratio $\Ac_{5156}(n)/\Ac_{5156}(n-1)$, not $b(n+1)/b(n)$. The statement as written is false
at $n=1$: the left-hand side is $\binom42/(2\binom32) = 1$ while $b(2)/b(1) = 3$. With the
intended sequence and shift the numerical agreement is exact:
$1, 3, 26/3, 323/13, \dots$ on both sides for $n = 1,2,3,4$. We record also that, granted
Kuperberg's product formula for A005156, the OEIS conjecture reduces to a factorial identity, so
that its substance is Kuperberg's theorem rather than the displayed ratio. We have not
formalized a corrected version.

\item The file \texttt{OEIS/100434.lean} defines $b(n) = c(n+1)$ for even $n$, omitting
a minus sign. All three of its conjectures are then false at $n=0$, where $c(0)+d(0) = 3$ but
$c(1) = -3$. Section~\ref{sec:a100434} discusses the intended reading and proves the three
identities in the corrected form.
\end{itemize}

Consistent with a policy of reporting rather than silently patching third-party corpora, we have
not modified those files; our corrected A100434 development lives in a separate namespace.

\section{The Kimberling floor formula (A108211)}
\label{sec:a108211}

\subsection{Statement and notation}

\begin{theorem}[Kimberling's conjecture]
\label{thm:a108211}
For every integer $n \ge 1$,
\[
  \left\lfloor \frac{1}{D(n)} \right\rfloor = 16n^2 + 1,
  \qquad\text{where}\qquad
  D(n) \;:=\; \frac{1}{4n} - \log 2 + \sum_{k=n+1}^{2n} \frac1k .
\]
\end{theorem}

Throughout this section we work with the following functions of a real variable $x \ge 1$. Put
\begin{align}
  A(x) &:= 1024x^5 + 5888x^4 + 12928x^3 + 13312x^2 + 6148x + 843,
  \label{eq:Apoly}\\
  B(x) &:= 4\,(4x+1)(2x+1)(4x+3)(4x+5)(2x+3)(4x+7),
  \label{eq:Bpoly}\\
  h(x) &:= \frac{A(x)}{B(x)}, \qquad
  g(x) := \frac{60}{(4x+1)^7},
  \label{eq:hg}
\end{align}
and
\begin{equation}
  f_1(x) := \frac{1}{4x} - \frac{1}{16x^2+1},
  \qquad
  f_2(x) := \frac{1}{4x} - \frac{1}{16x^2+2}.
  \label{eq:f1f2}
\end{equation}
Both $A$ and $B$ have positive coefficients, so $h(x) > 0$ and $g(x) > 0$ for $x \ge 0$.
Finally, for an integer $n \ge 1$ set
\begin{equation}
  T(n) \;:=\; \sum_{m \ge n} \frac{1}{(2m+1)(2m+2)} .
  \label{eq:Tdef}
\end{equation}
The series converges, being dominated by $\sum m^{-2}$.

The shape of the proof is this. Step~1 identifies $D(n) = \frac{1}{4n} - T(n)$. Steps~2 and~3
sandwich $T(n)$ between $h(n)$ and $h(n)+g(n)$ by a telescoping certificate. Step~4 shows that
this sandwich lies strictly inside $(f_1(n), f_2(n))$, which is exactly what pins the floor.
The reader should note the margin at stake: the interval
$\bigl(f_1(n), f_2(n)\bigr)$ has width
\[
  f_2(n) - f_1(n) = \frac{1}{16n^2+1} - \frac{1}{16n^2+2}
    = \frac{1}{(16n^2+1)(16n^2+2)} \;=\; \Theta(n^{-4}),
\]
while the sandwich supplied by Step~3 has width $g(n) = 60/(4n+1)^7 = \Theta(n^{-7})$. The
whole design of the certificate $h$ --- a ratio of a quintic by a sextic --- is dictated by the
need to beat $n^{-4}$; a cruder telescoping certificate produces an uncertainty of the same
order as the window and proves nothing.

\subsection{Step 1: reduction to the tail of the paired alternating harmonic series}

\begin{lemma}[Catalan's identity]
\label{lem:catalan}
For every integer $n \ge 0$, $\displaystyle \sum_{k=n+1}^{2n} \frac1k
= \sum_{j=1}^{2n} \frac{(-1)^{j+1}}{j}$.
\end{lemma}

\begin{proof}
Writing $H_N = \sum_{j \le N} 1/j$, the left side is $H_{2n} - H_n$. Since
$H_n = 2\sum_{i=1}^{n} \frac{1}{2i}$ is twice the sum of the reciprocals of the even integers up
to $2n$,
\[
  H_{2n} - H_n \;=\; \sum_{j=1}^{2n} \frac1j - 2\!\!\sum_{\substack{j \le 2n \\ j \text{ even}}}\!\! \frac1j
  \;=\; \sum_{j=1}^{2n} \frac{(-1)^{j+1}}{j}. \qedhere
\]
\end{proof}

\begin{lemma}[Mercator's series at $x=1$]
\label{lem:mercator}
$\displaystyle \log 2 = \sum_{j=1}^{\infty} \frac{(-1)^{j+1}}{j}$, and for every $N \ge 0$,
$\bigl| \log 2 - \sum_{j=1}^{N} (-1)^{j+1}/j \bigr| \le \frac{1}{N+1}$.
\end{lemma}

\begin{proof}
For $x \in [0,1]$ and $N \ge 0$, summing the finite geometric series gives the exact identity
\[
  \frac{1}{1+x} \;=\; \sum_{j=1}^{N} (-1)^{j+1} x^{j-1} \;+\; \frac{(-1)^N x^N}{1+x}.
\]
Integrating over $[0,1]$ yields
$\log 2 = \sum_{j=1}^{N} (-1)^{j+1}/j + (-1)^N \int_0^1 \frac{x^N}{1+x}\,dx$, and
$0 \le \int_0^1 \frac{x^N}{1+x}\,dx \le \int_0^1 x^N dx = \frac{1}{N+1}$.
\end{proof}

\begin{lemma}
\label{lem:step1}
For every integer $n \ge 1$, $\;D(n) = \dfrac{1}{4n} - T(n)$.
\end{lemma}

\begin{proof}
Let $S_N := \sum_{j=1}^{N} (-1)^{j+1}/j$, so $S_N \to \log 2$ by Lemma~\ref{lem:mercator}. By
Lemma~\ref{lem:catalan}, $D(n) = \frac{1}{4n} - \bigl(\log 2 - S_{2n}\bigr)$, so it suffices to
prove $\log 2 - S_{2n} = T(n)$.

Consider the partial sums of the paired series \eqref{eq:Tdef}. For $M \ge 0$,
\[
  \sum_{m=n}^{n+M-1} \frac{1}{(2m+1)(2m+2)}
  \;=\; \sum_{m=n}^{n+M-1} \left( \frac{1}{2m+1} - \frac{1}{2m+2} \right)
  \;=\; S_{2n+2M} - S_{2n},
\]
because the $M$ consecutive pairs of terms $\bigl(+\frac{1}{2m+1}, -\frac{1}{2m+2}\bigr)$,
$m = n, \dots, n+M-1$, are precisely the terms of index $2n+1, 2n+2, \dots, 2n+2M$ of the
alternating series, in order. Thus the $M$-th partial sum of the paired series equals
$S_{2n+2M} - S_{2n}$: the partial sums of the paired series form (a shift of) a
\emph{subsequence} of the partial sums of the alternating series, namely the even-indexed one.
A subsequence of a convergent sequence converges to the same limit, so letting $M \to \infty$
gives $T(n) = \log 2 - S_{2n}$.
\end{proof}

\begin{remark}
This is the point at which the review objected to the phrase ``grouping the tail in consecutive
pairs is legitimate since the partial sums converge''. Consecutive grouping of a convergent
series does preserve the sum, but the reason is the subsequence argument just given, not
absolute convergence --- which the alternating harmonic series does not have. In the Lean
development the issue disappears entirely: the identity
$\sum_{m<n} \frac{1}{(2m+1)(2m+2)} = H_{2n} - H_n$ is proved there by a finite induction on $n$,
so no regrouping of an infinite series is ever performed
(Section~\ref{sec:formal}).
\end{remark}

\subsection{Step 2: the telescoping certificate}

\begin{lemma}[Defect identity]
\label{lem:defect}
For every real $m \ge 0$,
\[
  \Delta(m) \;:=\; \frac{1}{(2m+1)(2m+2)} - \bigl( h(m) - h(m+1) \bigr)
  \;=\; \frac{63\,(16m^2 + 120m + 119)}{2\,P(m)},
\]
where
\[
  P(m) := (m+1)(2m+1)(2m+3)(2m+5)(4m+1)(4m+3)(4m+5)(4m+7)(4m+9)(4m+11).
\]
In particular $\Delta(m) > 0$ for every $m \ge 0$.
\end{lemma}

\begin{proof}
This is an identity between rational functions. Clearing denominators turns it into an identity
between polynomials of degree $12$, which is verified by expansion; the verification is
mechanical and was performed independently three times (by computer algebra over $\Q$, by an
adversarial reviewer recomputing it symbolically, and by Lean's \texttt{field\_simp; ring}
in the file \texttt{A108211.lean}, where the same identity appears as \texttt{delta\_eq} with
the numerator written $126(16m^2+120m+119)$ over $4P(m)$).

Positivity is immediate: for $m \ge 0$ every factor of $P(m)$ is positive and
$16m^2+120m+119 \ge 119 > 0$.
\end{proof}

\begin{lemma}[Majorant certificate]
\label{lem:majorant}
For every real $m \ge 1$,
\[
  0 \;<\; \Delta(m) \;<\; g(m) - g(m+1) \;=\; \frac{60}{(4m+1)^7} - \frac{60}{(4m+5)^7}.
\]
More precisely, with
\[
  D_2(m) := 2\,(m+1)(2m+1)(2m+3)(2m+5)(4m+1)^7(4m+3)(4m+5)^7(4m+7)(4m+9)(4m+11),
\]
which is positive for $m \ge 0$, one has $\;g(m) - g(m+1) - \Delta(m) = N_2(m)/D_2(m)$, where $N_2$ is a polynomial of degree
$14$ all of whose coefficients in the shifted variable $t = m-1$ --- that is, the coefficients of
$N_2(1+t)$ --- are the positive integers listed in Table~\ref{tab:C2}. Consequently
$N_2(m) > 0$ for every $m \ge 1$.
\end{lemma}

\begin{proof}
The displayed rational identity is verified by clearing denominators and expanding, exactly as
in Lemma~\ref{lem:defect}; note that $g(m+1) = 60/(4(m+1)+1)^7 = 60/(4m+5)^7$. Substituting
$m = 1+t$ with $t \ge 0$ turns $N_2(m)$ into $N_2(1+t) = \sum_{k=0}^{14} c_k t^k$ with every
$c_k > 0$ (Table~\ref{tab:C2}); the smallest coefficient is $c_{14} = 11{,}274{,}289{,}152$ and
the constant term is $c_0 = 1{,}646{,}809{,}468{,}266{,}375$. A sum of nonnegative multiples of
nonnegative quantities, with a positive constant term, is positive; hence $N_2(1+t) > 0$ for all
$t \ge 0$, i.e. $N_2(m) > 0$ for all $m \ge 1$. Since $D_2(m) > 0$ there, the claim follows.
The lower bound $\Delta(m) > 0$ is Lemma~\ref{lem:defect}.
\end{proof}

\begin{remark}
The device of exhibiting a univariate polynomial positivity certificate by shifting to
$x = 1+t$ and displaying nonnegative coefficients is what makes these lemmas checkable both by
hand and by a proof assistant: after the shift, positivity is a \emph{linear} consequence of the
facts $t^k \ge 0$, so no nonlinear arithmetic is needed. This is exploited in
Section~\ref{sec:formal}.
\end{remark}

\subsection{Step 3: the sandwich, with the first-term deficit retained}

\begin{lemma}[Decay certificate]
\label{lem:decay}
For every real $x \ge 1$, $\;0 < h(x) \le 1/x$. In particular $h(x) \to 0$ as $x \to \infty$.
\end{lemma}

\begin{proof}
Positivity is clear. The inequality $h(x) \le 1/x$ is equivalent to $x\,A(x) \le B(x)$, both
sides being polynomials of degree $6$ and $B(x) > 0$ for $x \ge 1$. Substituting $x = 1+t$,
\[
  B(1+t) - (1+t)A(1+t) = 3072t^6 + 37120t^5 + 184448t^4 + 482304t^3 + 699916t^2 + 534509t + 167757,
\]
all of whose coefficients are positive (Table~\ref{tab:C5}); hence the difference is positive
for $t \ge 0$.
\end{proof}

\begin{lemma}[Sandwich]
\label{lem:sandwich}
For every integer $n \ge 1$,
\[
  h(n) \;<\; T(n) \;\le\; h(n) + g(n) - \delta(n) \;<\; h(n) + g(n),
  \qquad \delta(n) := \frac{N_2(n)}{D_2(n)} > 0,
\]
with $N_2, D_2$ as in Lemma~\ref{lem:majorant}.
\end{lemma}

\begin{proof}
By Lemma~\ref{lem:defect}, for every integer $m \ge n$,
\begin{equation}
  \frac{1}{(2m+1)(2m+2)} \;=\; \bigl(h(m) - h(m+1)\bigr) + \Delta(m).
  \label{eq:split}
\end{equation}
Summing \eqref{eq:split} over $n \le m \le M$ and telescoping,
\begin{equation}
  \sum_{m=n}^{M} \frac{1}{(2m+1)(2m+2)}
  \;=\; h(n) - h(M+1) + \sum_{m=n}^{M} \Delta(m).
  \label{eq:partial}
\end{equation}
The terms $\Delta(m)$ are positive (Lemma~\ref{lem:defect}) and, by Lemma~\ref{lem:majorant},
bounded above by $g(m)-g(m+1)$, whose partial sums telescope to
$g(n) - g(M+1) \le g(n)$. Hence $\sum_{m \ge n} \Delta(m)$ converges, say to $\Sigma_n \ge 0$;
and $h(M+1) \to 0$ by Lemma~\ref{lem:decay}. Letting $M \to \infty$ in \eqref{eq:partial},
\begin{equation}
  T(n) \;=\; h(n) + \Sigma_n .
  \label{eq:Tsplit}
\end{equation}

\emph{Lower bound.} $\Sigma_n \ge \Delta(n) > 0$, so $T(n) > h(n)$.

\emph{Upper bound, with the deficit retained.} Set
$\delta(m) := \bigl(g(m) - g(m+1)\bigr) - \Delta(m) = N_2(m)/D_2(m)$, which is
\emph{strictly} positive for every $m \ge 1$ by Lemma~\ref{lem:majorant}. For any $M \ge n$,
\[
  \sum_{m=n}^{M} \Delta(m)
  = \sum_{m=n}^{M} \bigl(g(m)-g(m+1)\bigr) - \sum_{m=n}^{M} \delta(m)
  = g(n) - g(M+1) - \sum_{m=n}^{M} \delta(m)
  \;\le\; g(n) - \delta(n),
\]
where the last step drops $-g(M+1) \le 0$ and all terms $\delta(m) \ge 0$ with $m > n$, keeping
only the term $m=n$. The bound $g(n) - \delta(n)$ does not depend on $M$, so letting
$M \to \infty$ gives $\Sigma_n \le g(n) - \delta(n) < g(n)$, and \eqref{eq:Tsplit} finishes the
proof.
\end{proof}

\begin{remark}
The last paragraph is deliberately fussy. Term-wise strict inequalities between infinitely many
pairs of terms do \emph{not} imply a strict inequality between the sums --- the deficits may
tend to $0$ fast enough for the sums to agree. Retaining the single deficit $\delta(n)$, an
explicit positive rational, before passing to the limit removes the difficulty. (For the
purposes of Theorem~\ref{thm:a108211} the non-strict bound $T(n) \le h(n)+g(n)$ would in fact
suffice, because Lemma~\ref{lem:window} is strict; this is the route taken in the Lean file.
We give the strict version because it is what the informal statement asserts.)
\end{remark}

\subsection{Step 4: the window inequalities and the proof of Theorem~\ref{thm:a108211}}

\begin{lemma}[Window]
\label{lem:window}
For every real $x \ge 1$,
\[
  f_1(x) \;<\; h(x)
  \qquad\text{and}\qquad
  h(x) + g(x) \;<\; f_2(x).
\]
\end{lemma}

\begin{proof}
\emph{Lower window.} Since $f_1(x) = \dfrac{16x^2-4x+1}{4x(16x^2+1)}$ and $h(x) = A(x)/B(x)$
with both denominators positive for $x \ge 1$, the inequality $f_1(x) < h(x)$ is equivalent to
$N_3(x) > 0$, where
\[
  N_3(x) := A(x)\cdot 4x(16x^2+1) - (16x^2-4x+1)\cdot B(x),
\]
a polynomial of degree $5$ once the degree-$8$ terms cancel. Explicitly, after the substitution
$x = 1+t$,
\[
  N_3(1+t) = 256t^5 + 2816t^4 + 12096t^3 + 24688t^2 + 22223t + 6756,
\]
all of whose coefficients are positive (Table~\ref{tab:C3}); hence $N_3(1+t) > 0$ for $t \ge 0$.
Equivalently, over the positive common denominator
$D_3(x) = 4x(2x+1)(2x+3)(4x+1)(4x+3)(4x+5)(4x+7)(16x^2+1)$ one has
$h(x)-f_1(x) = N_3(x)/D_3(x) > 0$; at $x=1$ the exact value is
$h(1)-f_1(1) = \tfrac{563}{294525} > 0$.

\emph{Upper window.} Since $f_2(x) = \dfrac{8x^2-2x+1}{4x(8x^2+1)}$ and
\[
  h(x) + g(x)
  = \frac{A(x)(4x+1)^6 + 240\,(2x+1)(4x+3)(4x+5)(2x+3)(4x+7)}
         {4(4x+1)^7 (2x+1)(4x+3)(4x+5)(2x+3)(4x+7)},
\]
the inequality $h(x)+g(x) < f_2(x)$ is equivalent to $N_4(x) > 0$ where $N_4$ is the numerator
of $f_2 - h - g$ over the positive common denominator
$D_4(x) = 4x(2x+1)(2x+3)(4x+1)^7(4x+3)(4x+5)(4x+7)(8x^2+1)$. Substituting $x = 1+t$,
\begin{align*}
  N_4(1+t) = \;& 393216\,t^9 + 6291456\,t^8 + 43757568\,t^7 + 169863168\,t^6 + 404609280\,t^5 \\
  &+ 614989440\,t^4 + 598254960\,t^3 + 359583120\,t^2 + 120914895\,t + 17203050,
\end{align*}
all of whose coefficients are positive (Table~\ref{tab:C4}); hence $N_4(1+t) > 0$ for $t \ge 0$.
At $x=1$ the exact value of the difference is
\[
  f_2(1) - h(1) - g(1) \;=\; \frac{12743}{21656250} \;>\; 0 .
\]
(This exact rational replaces the decimal ``$0.000588\ldots$'' of an earlier draft; no decimal
approximation occurs anywhere in the present argument.)
\end{proof}

\begin{proof}[Proof of Theorem~\ref{thm:a108211}]
Fix an integer $n \ge 1$. Chaining Lemmas~\ref{lem:sandwich} and~\ref{lem:window} at $x = n$,
\[
  f_1(n) \;<\; h(n) \;<\; T(n) \;\le\; h(n) + g(n) - \delta(n) \;<\; h(n) + g(n) \;<\; f_2(n),
\]
so $f_1(n) < T(n) < f_2(n)$. By Lemma~\ref{lem:step1}, $D(n) = \frac{1}{4n} - T(n)$, and
subtracting the chain from $\frac{1}{4n}$ reverses it:
\[
  \frac{1}{4n} - f_2(n) \;<\; D(n) \;<\; \frac{1}{4n} - f_1(n),
  \qquad\text{that is,}\qquad
  \frac{1}{16n^2+2} \;<\; D(n) \;<\; \frac{1}{16n^2+1}
\]
by the definitions \eqref{eq:f1f2}. In particular $D(n) > \frac{1}{16n^2+2} > 0$, so we may
invert; inversion of positive reals reverses strict inequalities, giving
\[
  16n^2 + 1 \;<\; \frac{1}{D(n)} \;<\; 16n^2 + 2 .
\]
Since $16n^2+1$ is an integer and $1/D(n)$ lies strictly between it and its successor,
$\lfloor 1/D(n) \rfloor = 16n^2+1$.
\end{proof}

\begin{remark}
The positivity of $D(n)$, established here \emph{before} inversion, is not a formality. By
Lemma~\ref{lem:step1} the assertion $D(n) > 0$ says exactly that $T(n) < \frac{1}{4n}$, i.e.
that $\sum_{k=n+1}^{2n} \frac1k > \log 2 - \frac{1}{4n}$, and since
$T(n) = \frac{1}{4n} - \frac{1}{16n^2} + O(n^{-4})$ this is a genuine inequality, delivered here
by the upper window bound. Note also that no appeal to the irrationality of $\log 2$ is made
anywhere: every inequality above is strict by certificate, so the possibility of $1/D(n)$ landing
exactly on an integer is excluded outright rather than by an arithmetic argument.
\end{remark}

\section{The asymptotics of A114831}
\label{sec:a114831}

\begin{theorem}
\label{thm:a114831}
Define $a : \Z_{\ge 1} \to \Z_{\ge 1}$ by $a(1) = 1$, $a(2) = 2$ and
\[
  a(n) \;=\; a(n-1) + \left\lfloor \frac{2\,a(n-1)\,a(n-2)}{a(n-1)+a(n-2)} \right\rfloor
  \qquad (n \ge 3),
\]
so that $a = 1, 2, 3, 5, 8, 14, 24, 41, 71, 122, 211, \dots$. Then
\[
  \lim_{n \to \infty} \frac{a(n+1)}{a(n)} \;=\; \sqrt3 ,
\]
and for every $n \ge 4$ one has the explicit bound
\begin{equation}
  \left| \frac{a(n)}{a(n-1)} - \sqrt3 \right|
  \;\le\; (2-\sqrt3)\,2^{-(n-2)} + \frac{4}{n} + 2^{-(n-2)/2}.
  \label{eq:a114831bound}
\end{equation}
\end{theorem}

Write $A_n := a(n)$ and, for $n \ge 2$, $R_n := A_n/A_{n-1}$.

\begin{proof}
\emph{Linear growth.} For positive integers $x, y$,
\[
  2xy - x - y \;=\; x(y-1) + y(x-1) \;\ge\; 0,
\]
so $2xy/(x+y) \ge 1$ and therefore the integer $\lfloor 2xy/(x+y) \rfloor$ is at least $1$.
Hence $A_n \ge A_{n-1} + 1$ for $n \ge 3$, and since $A_1 = 1$, $A_2 = 2$,
\begin{equation}
  A_n \;\ge\; n \qquad (n \ge 1).
  \label{eq:linear}
\end{equation}
In particular $A_n > A_{n-1} \ge 1$, so $R_n > 1$ for every $n \ge 2$.

\emph{The exact ratio recursion.} For $n \ge 3$ let
\[
  \theta_n \;:=\; \frac{2A_{n-1}A_{n-2}}{A_{n-1}+A_{n-2}}
    - \left\lfloor \frac{2A_{n-1}A_{n-2}}{A_{n-1}+A_{n-2}} \right\rfloor \;\in\; [0,1)
\]
be the fractional defect of the floor. Since $R_{n-1} = A_{n-1}/A_{n-2}$,
\[
  \frac{2A_{n-1}A_{n-2}}{A_{n-1}+A_{n-2}}
  \;=\; \frac{2A_{n-1}}{\;(A_{n-1}+A_{n-2})/A_{n-2}\;}
  \;=\; \frac{2A_{n-1}}{R_{n-1}+1},
\]
so dividing the defining recurrence $A_n = A_{n-1} + \frac{2A_{n-1}A_{n-2}}{A_{n-1}+A_{n-2}} - \theta_n$
by $A_{n-1}$ gives the exact identity
\begin{equation}
  R_n \;=\; 1 + \frac{2}{R_{n-1}+1} - \frac{\theta_n}{A_{n-1}}
  \;=\; F(R_{n-1}) - \frac{\theta_n}{A_{n-1}},
  \qquad F(x) := 1 + \frac{2}{x+1}.
  \label{eq:ratiorec}
\end{equation}

\emph{Fixed point and contraction.} Put $s := \sqrt3$. Since $(s-1)(s+1) = s^2-1 = 2$, we have
$2/(s+1) = s-1$ and therefore
\begin{equation}
  F(s) \;=\; 1 + (s-1) \;=\; s.
  \label{eq:fixed}
\end{equation}
Moreover, for all $x, y \ge 1$,
\begin{equation}
  |F(x) - F(y)| \;=\; \frac{2\,|x-y|}{(x+1)(y+1)} \;\le\; \frac{|x-y|}{2},
  \label{eq:lipschitz}
\end{equation}
because $(x+1)(y+1) \ge 4$. Thus $F$ is a $\tfrac12$-contraction on $[1,\infty)$, a set that
contains both $s = 1.732\ldots$ and every $R_m$ with $m \ge 2$.

\emph{Error recursion.} Let $E_n := |R_n - s|$. Combining \eqref{eq:ratiorec},
\eqref{eq:fixed}, \eqref{eq:lipschitz}, $R_{n-1} \ge 1$, $s \ge 1$, $0 \le \theta_n < 1$, and
\eqref{eq:linear}, we obtain for every $n \ge 3$
\begin{equation}
  E_n \;\le\; |F(R_{n-1}) - F(s)| + \frac{\theta_n}{A_{n-1}}
  \;\le\; \tfrac12 E_{n-1} + \frac{1}{A_{n-1}}
  \;\le\; \tfrac12 E_{n-1} + \frac{1}{n-1}.
  \label{eq:erec}
\end{equation}
Since $R_2 = A_2/A_1 = 2$, we have $E_2 = 2 - \sqrt3$, and iterating \eqref{eq:erec} gives
\begin{equation}
  E_n \;\le\; 2^{-(n-2)}(2-\sqrt3) + \sum_{j=3}^{n} \frac{2^{-(n-j)}}{j-1}.
  \label{eq:iterated}
\end{equation}

\emph{Explicit tail bound.} Let $n \ge 4$ and put $K := \lfloor (n-2)/2 \rfloor$. Substituting
$k = n-j$,
\[
  \sum_{j=3}^{n} \frac{2^{-(n-j)}}{j-1} \;=\; \sum_{k=0}^{n-3} \frac{2^{-k}}{n-k-1}.
\]
For $0 \le k \le K$ we have $n-k-1 \ge n-K-1 \ge n/2$ (because $K \le (n-2)/2$), while for
$K+1 \le k \le n-3$ we have $n-k-1 \ge 2$. Therefore
\begin{equation}
  \sum_{k=0}^{n-3} \frac{2^{-k}}{n-k-1}
  \;\le\; \frac2n \sum_{k=0}^{K} 2^{-k} + \frac12 \sum_{k=K+1}^{\infty} 2^{-k}
  \;\le\; \frac4n + 2^{-K-1}
  \;\le\; \frac4n + 2^{-(n-2)/2},
  \label{eq:tail}
\end{equation}
using $\sum_{k \ge 0} 2^{-k} = 2$ and $K+1 \ge (n-2)/2$. Combining \eqref{eq:iterated} and
\eqref{eq:tail} gives exactly \eqref{eq:a114831bound}. Each of the three terms on the right of
\eqref{eq:a114831bound} tends to $0$, so $E_n \to 0$, i.e. $R_n \to \sqrt3$. Replacing $n$ by
$n+1$ gives $\lim_{n\to\infty} a(n+1)/a(n) = \sqrt3$.
\end{proof}

\begin{remark}[Two independent derivations]
Theorem~\ref{thm:a114831} was proved twice, independently and without communication between the
two solvers. The first derivation (Claude, Fable 5) established an invariant interval
$R_n \in [1.59, 2]$ for $n$ large, used the exact algebraic contraction identity
\[
  F(r) - \sqrt3 \;=\; \frac{(r-\sqrt3)(1-\sqrt3)}{1+r},
\]
which gives a contraction factor $(\sqrt3-1)/(1+r) \le 0.283$ on that interval, and controlled
the forcing term by the Fibonacci lower bound $A_n \ge F_n$ (valid because
$\lfloor 2xy/(x+y) \rfloor \ge \min(x,y)$, so $A_n \ge A_{n-1}+A_{n-2}$). The second derivation,
reproduced above, was obtained by OpenAI GPT-5.6 and is strictly simpler: the Lipschitz bound
\eqref{eq:lipschitz} is global on $[1,\infty)$, so no invariant interval is needed, and the
crude linear bound \eqref{eq:linear} suffices in place of the Fibonacci one. Cross-checking the
two line by line --- both reach $\sqrt3$ through the same fixed-point equation but with
different contraction constants ($0.283$ versus $1/2$) and different growth inputs --- is the
reason we adopted the second for formalization. The stronger Fibonacci growth would improve
\eqref{eq:a114831bound} from $O(1/n)$ to a geometric rate; we did not pursue this, since the
conjecture asks only for the limit.
\end{remark}

\begin{remark}[Prior work]
While this paper was in preparation, Kenta Kitamura (GitHub user \texttt{KitaKen1}) independently
proved OEIS A114831's \texttt{conjecture3} --- the same statement,
$a(n+1)/a(n) \to \sqrt3$ --- with a kernel-checked Lean~4 proof
(repository \url{https://github.com/KitaKen1/oeis-a114831-asymptotic}), submitted as
\texttt{google-deepmind/formal-conjectures} pull request \#4969, opened 15~August~2026 and merged
16~August~2026, which marked the conjecture \texttt{research solved} upstream. That submission
predates our own derivation and formalization, completed the same day; we became aware of it only
after our proof was finished. We therefore regard Theorem~\ref{thm:a114831} as an independent,
essentially simultaneous proof, obtained by a different argument --- the global
$\tfrac12$-Lipschitz contraction \eqref{eq:lipschitz} used above, rather than Kitamura's method ---
and we make no claim to priority.
\end{remark}

\section{The Dement identities (A100434)}
\label{sec:a100434}

\subsection{The sequences and the sign convention}

A100434 is the expansion of $(1+x)(3+x)/(1+6x^2+x^4)$, i.e. the integer sequence
\begin{equation}
  a(0)=3,\quad a(1)=4,\quad a(2)=-17,\quad a(3)=-24,
  \qquad a(n+4) = -6a(n+2) - a(n),
  \label{eq:a100434}
\end{equation}
beginning $3, 4, -17, -24, 99, 140, -577, -816, 3363, \ldots$. Dement's comment introduces two
further solutions of the same recurrence,
\begin{align}
  c &: \quad c(0)=1,\; c(1)=-3,\; c(2)=-7,\; c(3)=17, \label{eq:cdef}\\
  d &: \quad d(0)=2,\; d(1)=4,\; d(2)=-10,\; d(3)=-24, \label{eq:ddef}
\end{align}
each extended by $x(n+4) = -6x(n+2)-x(n)$, and four derived sequences
\begin{equation}
  \begin{aligned}
    e(2k) &= \tfrac12 d(2k), & e(2k+1) &= -\tfrac12 d(2k),\\
    f(2k) &= f(2k+1) = \tfrac12 d(2k+1), & &\\
    g(2k) &= 0, & g(2k+1) &= c(2k+1),
  \end{aligned}
  \label{eq:efg}
\end{equation}
together with a sequence $b$ described as ``$c$ with even- and odd-indexed terms reversed''. The
conjecture is
\[
  c(n) + d(n) \;=\; e(n) + f(n) \;=\; g(n) + a(n) \;=\; b(n) \qquad (n \ge 0).
\]

The comment writes this reversal as $b(2k) = c(2k+1)$, $b(2k+1) = c(2k)$. Taken literally the
conjecture is false, and at once: $c(0) + d(0) = 1 + 2 = 3$ while $c(1) = -3$. The intended
reading is however unambiguous from the surrounding text, which introduces $c$ as ``the sequence
A001333, apart from signs'' and $d$ as ``A052542, apart from signs'', and which summarises the
three identities as saying that all three sums ``represent the sequence $c$ with even- and
odd-indexed terms reversed''. The sign pattern of $c$ is $+,-,-,+,+,-,-,+,\dots$ with period
four; reversing terms within consecutive pairs necessarily disturbs it, so a statement about the
reversal is a statement about $|c|$, i.e. one made ``apart from signs''. Restoring the sign so
that the identities hold requires exactly one change,
\begin{equation}
  b(2k) := -c(2k+1), \qquad b(2k+1) := c(2k),
  \label{eq:bdef}
\end{equation}
and this choice is independently corroborated: $-c(2k+1) = a(2k)$ is a known identity for this
family (it is stated in the OEIS comment itself, as ``$a(2n) = -c(2n+1)$''), so \eqref{eq:bdef}
makes $b$ the interleaving of the even part of $a$ with the even part of $c$,
\[
  b = 3,\; 1,\; -17,\; -7,\; 99,\; 41,\; -577,\; -239,\; 3363,\; 1393,\; \ldots
\]
which is $|c| = 1, 3, 7, 17, 41, 99, \ldots$ with consecutive terms transposed, carrying the sign
pattern of $a$. With \eqref{eq:bdef} the three identities hold for every $n$; they were verified
numerically for $n < 600$ before any proof was attempted, and are proved below.

\subsection{The identities}

\begin{theorem}
\label{thm:a100434}
Let $a, c, d$ be given by \eqref{eq:a100434}--\eqref{eq:ddef}, let $e, f, g$ be given by
\eqref{eq:efg}, and let $b$ be given by \eqref{eq:bdef}. Then for every integer $n \ge 0$,
\[
  c(n) + d(n) \;=\; b(n),
  \qquad
  e(n) + f(n) \;=\; b(n),
  \qquad
  g(n) + a(n) \;=\; b(n).
\]
\end{theorem}

\begin{proof}
All of $a, c, d$ satisfy the same order-four recurrence $x(n+4) = -6x(n+2) - x(n)$, which
couples indices of equal parity only. Consequently every statement below splits into an even and
an odd case, and each case is proved by two-step induction on $k$ (base cases $k = 0, 1$,
inductive step from $k, k+1$ to $k+2$), the step being a linear identity among the four values
involved. We record the six auxiliary facts; each is proved by exactly this scheme, and the base
cases are the displayed initial values.

\begin{align}
  c(2k) + d(2k) &= -c(2k+1), \label{eq:aux1}\\
  c(2k+1) + d(2k+1) &= c(2k), \label{eq:aux2}\\
  a(2k) &= -c(2k+1), \qquad a(2k+1) = d(2k+1), \label{eq:aux3}\\
  d(2k) &\equiv 0, \quad d(2k+1) \equiv 0 \pmod 2, \label{eq:aux4}\\
  d(2k) + d(2k+1) &= -2\,c(2k+1), \label{eq:aux5}\\
  d(2k+1) - d(2k) &= 2\,c(2k). \label{eq:aux6}
\end{align}

For instance, \eqref{eq:aux1} at $k=0$ reads $c(0)+d(0) = 1+2 = 3 = -c(1)$, and at $k=1$ it reads
$c(2)+d(2) = -7 + (-10) = -17 = -c(3)$; the inductive step is
\begin{align*}
  c(2k+4) + d(2k+4)
  &= -6\bigl(c(2k+2)+d(2k+2)\bigr) - \bigl(c(2k)+d(2k)\bigr)\\
  &= -6\bigl(-c(2k+3)\bigr) - \bigl(-c(2k+1)\bigr) \;=\; -c(2k+5),
\end{align*}
using the induction hypotheses at $k+1$ and $k$ and then the recurrence for $c$. The remaining
five are identical in form; \eqref{eq:aux4} is the same induction carried out for divisibility,
with base cases $d(0) = 2$, $d(2) = -10$, $d(1) = 4$, $d(3) = -24$.

\emph{First identity.} If $n = 2k$, then $c(n)+d(n) = -c(2k+1) = b(2k)$ by \eqref{eq:aux1} and
\eqref{eq:bdef}. If $n = 2k+1$, then $c(n)+d(n) = c(2k) = b(2k+1)$ by \eqref{eq:aux2} and
\eqref{eq:bdef}.

\emph{Second identity.} By \eqref{eq:aux4} the halvings in \eqref{eq:efg} are exact. If
$n = 2k$, then by \eqref{eq:efg} and \eqref{eq:aux5},
\[
  e(2k) + f(2k) = \tfrac12 d(2k) + \tfrac12 d(2k+1) = \tfrac12\bigl(d(2k)+d(2k+1)\bigr)
  = -c(2k+1) = b(2k).
\]
If $n = 2k+1$, then $f(2k+1) = \tfrac12 d(2k+1)$ and $e(2k+1) = -\tfrac12 d(2k)$, so by
\eqref{eq:aux6},
\[
  e(2k+1) + f(2k+1) = \tfrac12\bigl(d(2k+1) - d(2k)\bigr) = c(2k) = b(2k+1).
\]

\emph{Third identity.} If $n = 2k$, then $g(2k) = 0$ and $a(2k) = -c(2k+1) = b(2k)$ by
\eqref{eq:aux3}. If $n = 2k+1$, then $g(2k+1) + a(2k+1) = c(2k+1) + d(2k+1) = c(2k) = b(2k+1)$
by \eqref{eq:aux3} and \eqref{eq:aux2}.
\end{proof}

\begin{remark}
The identities are, as OEIS's own \texttt{easy} keyword suggests, not deep: they are six parallel
two-step inductions. The interest of the case lies elsewhere. First, the conjecture as it stands
in a formal corpus was refutable at the very first index, and this was found by evaluating the
formal statement numerically before attempting a proof --- a cheap check that we now run
routinely. Second, deciding \emph{which} corrected statement is ``the intended one'' is a
judgement about a natural-language source, not a mathematical deduction; we have set out the
evidence for \eqref{eq:bdef} above and flag it as the one interpretive step in this paper.
\end{remark}

\section{The Ordowski tanh-sum asymptotic (A114362)}
\label{sec:a114362}

\subsection{Statement and the Euler-product reduction}

\begin{theorem}[Ordowski's conjecture]
\label{thm:a114362}
For an integer $n \ge 2$ let
\[
  t(n) := \frac{\zeta(2n)}{\zeta(n)^2}, \qquad
  y(n) := \frac{1-t(n)}{1+t(n)}, \qquad
  S(n) := \frac{1}{2^n}+\frac{1}{3^n}+\frac{1}{5^n}+\frac{1}{7^n}.
\]
Then
\[
  y(n) \;=\; S(n) + O(11^{-n}), \qquad\text{explicitly}\qquad
  \bigl| y(n) - S(n) \bigr| \;\le\; 24 \cdot 11^{-n} \quad (n \ge 2).
\]
\end{theorem}

For a prime $p$ and $n \ge 2$ write $x_p := p^{-n} \in (0, \tfrac14]$ (the bound holding since
$p \ge 2$, $n \ge 2$) and
\[
  q_p := \frac{1-x_p}{1+x_p} \in (0,1).
\]

\begin{lemma}[Euler-product form]
\label{lem:ordowski-euler}
For every integer $n \ge 2$, $\displaystyle t(n) = \prod_p q_p$, the product being absolutely
convergent.
\end{lemma}

\begin{proof}
For $s>1$, $\zeta(s) = \prod_p (1-p^{-s})^{-1}$, the Euler product converging absolutely because
$\sum_p p^{-s} \le \sum_{m\ge2} m^{-s} < \infty$ and $-\log(1-u) \le \tfrac43 u$ for $0 \le u \le
\tfrac14$ bounds $\sum_p |\log(1-p^{-s})|$. Applying this at $s=n$ and $s=2n$,
\[
  t(n) \;=\; \frac{\zeta(2n)}{\zeta(n)^2}
  \;=\; \prod_p \frac{(1-p^{-n})^2}{1-p^{-2n}}
  \;=\; \prod_p \frac{1-p^{-n}}{1+p^{-n}} \;=\; \prod_p q_p.
\]
Absolute convergence of the last product follows from $|\log q_p| \le x_p + \tfrac43 x_p =
\tfrac73 x_p$ (using $\log(1+u)\le u$ for both factors) and $\sum_p x_p < \infty$.
\end{proof}

\subsection{Peeling off the primes $2,3,5,7$}

\begin{definition}
For $a,b \in [0,1)$ put $a \oplus b := \dfrac{a+b}{1+ab}$, and for $u \in (0,1]$ put
$\Phi(u) := \dfrac{1-u}{1+u}$.
\end{definition}

The operation $\oplus$ is the familiar relativistic velocity addition; it is commutative and
associative, and maps $[0,1)^2$ into $[0,1)$ since $1-(a\oplus b) = (1-a)(1-b)/(1+ab) > 0$. It is
related to $\Phi$ by the exact identity, valid for $a \in [0,1)$ and $R \in (0,1]$,
\begin{equation}
  \Phi\bigl(q(a)\,R\bigr) \;=\; a \oplus \Phi(R), \qquad q(a) := \frac{1-a}{1+a},
  \label{eq:peel}
\end{equation}
checked by clearing denominators on both sides: both equal
$\bigl((1+a)-(1-a)R\bigr) / \bigl((1+a)+(1-a)R\bigr)$.

\begin{lemma}[Peeling]
\label{lem:oplus-peel}
For $n \ge 2$ let $R_{11} := \prod_{p \ge 11} q_p \in (0,1)$ and $r := \Phi(R_{11})$, and put
$A := x_2 \oplus x_3 \oplus x_5 \oplus x_7$ (associativity makes the bracketing immaterial). Then
\[
  y(n) \;=\; A \oplus r.
\]
\end{lemma}

\begin{proof}
By Lemma~\ref{lem:ordowski-euler}, $t(n) = q_2 q_3 q_5 q_7 R_{11} = q(x_2)q(x_3)q(x_5)q(x_7)R_{11}$,
since $q_p = q(x_p)$ by definition. Applying \eqref{eq:peel} four times, innermost first,
\[
  y(n) = \Phi(t(n)) = x_2 \oplus \bigl(x_3 \oplus (x_5 \oplus (x_7 \oplus r))\bigr) = A \oplus r
\]
by associativity of $\oplus$.
\end{proof}

\subsection{The tail bound and the four-prime defect}

\begin{lemma}[Tail bound]
\label{lem:ordowski-tail}
For every integer $n \ge 2$, $0 \le r \le 24 \cdot 11^{-n}$.
\end{lemma}

\begin{proof}
For $u_j \in [0,1]$, $1 - \prod_{j\le m} u_j \le \sum_{j \le m}(1-u_j)$, by induction from
$1-uv = (1-u)+u(1-v) \le (1-u)+(1-v)$; passing to the limit over the (absolutely convergent) tail
product gives $1-R_{11} \le \sum_{p \ge 11}(1-q_p)$. Since $1-q_p = 2x_p/(1+x_p) \le 2x_p$,
\[
  1 - R_{11} \;\le\; 2\sum_{p \ge 11} p^{-n} \;\le\; 2\sum_{k=11}^{\infty} k^{-n}
  \;\le\; 2\cdot 11^{-n}\Bigl(1+\frac{11}{n-1}\Bigr) \;\le\; 24\cdot 11^{-n} \qquad (n \ge 2),
\]
the middle inequality by the integral comparison $\sum_{k\ge11}k^{-n} \le
11^{-n}+\int_{11}^\infty u^{-n}\,du$ and the last because $1+11/(n-1) \le 12$ for $n \ge 2$. Since
$R_{11} \ge 0$, $r = (1-R_{11})/(1+R_{11}) \le 1-R_{11} \le 24\cdot11^{-n}$; and $r\ge0$ since
$R_{11}\le1$.
\end{proof}

\begin{lemma}[Four-prime defect]
\label{lem:ordowski-defect}
Let $\delta(a,b) := a+b-(a\oplus b) = ab(a+b)/(1+ab) \ge 0$ for $a,b \ge 0$, and set
$b_5 := x_5\oplus x_7$, $b_3 := x_3\oplus b_5$, $F := S(n) - A$. Then
\[
  F \;=\; \delta(x_5,x_7) + \delta(x_3,b_5) + \delta(x_2,b_3) \;\le\; 20\cdot12^{-n} \qquad
  (n\ge2).
\]
\end{lemma}

\begin{proof}
Telescoping,
$F = (x_5+x_7-b_5)+(x_3+b_5-b_3)+(x_2+b_3-A) = \delta(x_5,x_7)+\delta(x_3,b_5)+\delta(x_2,b_3)$.
Since $0 < x_7 \le x_5 \le x_3 \le x_2$ (as $p \mapsto p^{-n}$ is decreasing), $b_5 \le x_5+x_7 \le
2x_5$ and $b_3 \le x_3+b_5 \le x_3+x_5+x_7 \le 3x_3$. Hence
\[
  \delta(x_5,x_7) \le x_5x_7(x_5+x_7) \le 2x_5^2x_7 = 2\cdot175^{-n},
\]
\[
  \delta(x_3,b_5) \le x_3b_5(x_3+b_5) \le x_3(2x_5)(3x_3) = 6x_3^2x_5 = 6\cdot45^{-n},
\]
and, since $x_2+b_3 \le x_2+3x_3 \le 4x_2$,
\[
  \delta(x_2,b_3) \le x_2b_3(x_2+b_3) \le x_2(3x_3)(4x_2) = 12x_2^2x_3 = 12\cdot12^{-n}.
\]
Since $45^{-n}, 175^{-n} \le 12^{-n}$, summing gives $F \le 12\cdot12^{-n}+6\cdot45^{-n}+
2\cdot175^{-n} \le 20\cdot12^{-n}$.
\end{proof}

\begin{proof}[Proof of Theorem~\ref{thm:a114362}]
For $a,b \ge 0$ with $a \in [0,1)$, $a\oplus b - a = b(1-a^2)/(1+ab) \ge 0$, and
$a\oplus b = (a+b)/(1+ab) \le a+b$; applying both to $y(n) = A\oplus r$
(Lemma~\ref{lem:oplus-peel}) gives $A \le y(n) \le A+r$. Since $S(n) = A+F$,
\[
  -F \;\le\; y(n)-S(n) \;\le\; r.
\]
By Lemma~\ref{lem:ordowski-defect}, $F \le 20\cdot12^{-n} \le 20\cdot11^{-n} \le 24\cdot11^{-n}$,
and by Lemma~\ref{lem:ordowski-tail}, $r \le 24\cdot11^{-n}$. Hence $|y(n)-S(n)| \le 24\cdot
11^{-n}$.
\end{proof}

\begin{remark}[Exact check at $n=2$, cross-checked three ways]
At $n=2$, $t(2) = \zeta(4)/\zeta(2)^2 = (\pi^4/90)/(\pi^2/6)^2 = 2/5$ (this is also
\texttt{OeisA114362.t\_two} in the Lean file), so $y(2) = 3/7$ --- matching the value $w(2)=3/7$
recorded independently in the \texttt{EXAMPLE} line of OEIS A348829, whose author's $w$ is the
same quantity as our $y$ under a different name. With $S(2) = \tfrac14+\tfrac19+\tfrac1{25}+
\tfrac1{49} = 18589/44100$, the exact defect is
\[
  y(2)-S(2) \;=\; \frac37 - \frac{18589}{44100} \;=\; \frac{311}{44100},
\]
positive and below $11^{-2}=1/121$ since $311\cdot121 = 37631 < 44100$. The reviewer of the
repaired proof recomputed every intermediate quantity at $n=2$ as an exact rational:
$A = 4669/11581$, $F = 9376309/510722100$, $R_{11} = 1625/1728$, $r = 103/3353$ --- values that
chain together correctly ($r=(1-R_{11})/(1+R_{11})$ recovers $103/3353$ from $R_{11}=1625/1728$)
and that satisfy every inequality of Lemmas~\ref{lem:ordowski-tail} and~\ref{lem:ordowski-defect}
at the exact value $n=2$, the case at which the first proof attempt had failed
(Section~\ref{sec:a114362-review} below).
\end{remark}

\begin{remark}[A sharper asymptotic, not formalized]
The repaired proof establishes, by the same method with the further prime $p=13$ peeled from
$R_{11}$, the stronger limit
\[
  \lim_{n\to\infty} \frac{y(n)-S(n)}{11^{-n}} \;=\; 1,
\]
which refines Theorem~\ref{thm:a114362} from a big-$O$ bound to an exact leading asymptotic. We
record it because it is part of the reviewed proof, but it is not part of the OEIS comment being
formalized, and we have not formalized it; Section~\ref{sec:formal} formalizes only the
$O(11^{-n})$ statement of Theorem~\ref{thm:a114362}, matching \texttt{formal-conjectures}'s
\texttt{conjecture2} verbatim.
\end{remark}

\begin{remark}[What is \emph{not} proved here]
Theorem~\ref{thm:a114362} leaves the sign of $y(n)-S(n)$ unresolved for general $n$; at $n=2$ it
is positive, but nothing above rules out the reverse. This is weaker than a related conjecture
Ordowski recorded the same day in OEIS A348829, restricted to even arguments:
$0 < w(2n) - S(2n) < 11^{-2n}$ for every $n>0$, a strict two-sided pin with no leading constant.
That statement, and Ordowski's further 2024 remark that the prime-zeta tail satisfies
$P(2n)-w(2n) \sim 12^{-2n}$ (consistent in order of magnitude with, though not implied by, the
$O(12^{-n})$ defect bound $F$ of Lemma~\ref{lem:ordowski-defect}), are neither proved nor
formalized here; we cite them in Section~\ref{sec:related} only as context for the same author's
closely related conjectures.
\end{remark}

\subsection{The review chain}
\label{sec:a114362-review}

\begin{remark}[Two refutations, two validations]
Theorem~\ref{thm:a114362} was proved, refuted, repaired, and re-validated in that order. A first
draft (Claude, Fable~5) reached the correct overall shape --- the Euler product, the
$\oplus$-peeling identity, and the $O(11^{-n})$ conclusion --- but asserted uniform bounds on the
intermediate $\oplus$-corrections, such as $x_5\cdot(x_7\oplus r) \le 2\cdot35^{-n}$, that are
simply false for small $n$. Two reviewers, working independently and without sight of each
other's report, located the same defect: Qwen3.8-Max showed the bound fails already at $n=2$ by
bounding the finite tail product over $p\in\{11,13,17,19,23\}$ below $23/24$, giving
$r > 1/47$ and hence $x_5\cdot(x_7\oplus r) > 2/1225 = 2\cdot35^{-2}$; GPT-5.6 (``Pro'' effort)
went further and exhibited the exact witness, inverting the $\oplus$-relations at $n=2$ to obtain
$x_3\oplus(x_5\oplus(x_7\oplus r)) = 1/5$, $x_5\oplus(x_7\oplus r) = 1/11$, and
$x_7\oplus r = 7/137$, so that $x_5\cdot(x_7\oplus r) = 7/3425$ exceeds the claimed bound
$2\cdot35^{-2}=2/1225$ by $69/167825$. Both reviewers returned \textsc{invalid}.

The repair --- Lemmas~\ref{lem:ordowski-tail} and~\ref{lem:ordowski-defect} above, which replace
the uniform pointwise bound on each nested correction with the single exact telescoping identity
$F = \delta(x_5,x_7)+\delta(x_3,b_5)+\delta(x_2,b_3)$ and a size comparison between the $x_p$
rather than an asymptotic estimate of each nested term --- was produced by GPT-5.6 (``Pro''
effort), the same model that had supplied the counterexample. Because the repair therefore did
not pass through an independent solver, we treat its own review as the operative gate: it was
checked, in separate sessions and blind to each other, by GPT-5.6 through \texttt{codex} and by
Qwen3.8-Max. Both returned \textsc{valid}; the \texttt{codex} review is the source of the exact
rationals reproduced in the remark above, obtained by recomputing $A$, $F$, $R_{11}$, and $r$
symbolically at $n=2$ rather than accepting the small-$n$ estimates on faith.
\end{remark}

\section{Formal verification}
\label{sec:formal}

\subsection{Files, environment, and axiom audits}

The four developments are self-contained Lean~4 files, each importing all of \texttt{mathlib}.
The toolchain is \texttt{leanprover/lean4:v4.34.0-rc1} with \texttt{mathlib} pinned at revision
\texttt{v4.34.0-rc1}. Table~\ref{tab:lean} summarises them. Elaboration times are wall-clock for
\texttt{lake env lean <file>} on an Apple-silicon laptop with a warm module cache and exclude
the one-off cost of building \texttt{mathlib}.

\begin{table}[ht]
\centering
\begin{tabular}{@{}llrr@{}}
\toprule
File & Main theorem & Lines & Elaboration \\
\midrule
\texttt{A108211.lean} & \texttt{a108211} & 509 & 5.6\,s \\
\texttt{A114831.lean} & \texttt{a114831} & 327 & 4.3\,s \\
\texttt{A100434.lean} & \texttt{conjecture1/2/3} & 282 & 2.4\,s \\
\texttt{A114362.lean} & \texttt{conjecture2} & 577 & 34.6\,s \\
\bottomrule
\end{tabular}
\caption{The Lean~4 developments. No file contains \texttt{sorry}, \texttt{native\_decide}, or
any \texttt{axiom} declaration.}
\label{tab:lean}
\end{table}

\texttt{A114362.lean} is markedly the slowest of the four to elaborate, consistent with its being
the only development that unfolds \texttt{mathlib}'s Euler-product theorem for the Riemann zeta
function; the figure above is the mean of three warm-cache runs (35.3\,s, 31.1\,s, 37.5\,s), and
the first, cold-cache invocation in a fresh shell took 61.9\,s.

Every file ends with an axiom audit. The complete output is:
\begin{quote}\ttfamily\small
\begin{tabbing}
'a108211' depends on axioms: [propext, Classical.choice, Quot.sound]\\
'a114831' depends on axioms: [propext, Classical.choice, Quot.sound]\\
'OeisA100434Corrected.conjecture1' depends on axioms: [propext, Classical.choice, Quot.sound]\\
'OeisA100434Corrected.conjecture2' depends on axioms: [propext, Classical.choice, Quot.sound]\\
'OeisA100434Corrected.conjecture3' depends on axioms: [propext, Classical.choice, Quot.sound]\\
'OeisA114362.conjecture2' depends on axioms: [propext, Classical.choice, Quot.sound]
\end{tabbing}
\end{quote}
These three are the standard axioms of Lean's classical foundation, on which \texttt{mathlib}
itself rests; in particular nothing is assumed beyond them, and no result is imported from an
unverified oracle. We note explicitly that \texttt{native\_decide}, which would place the
compiler's evaluator in the trusted base, is used nowhere.

\subsection{Statement fidelity}

\emph{A108211.} The Lean statement is
\begin{quote}\ttfamily\small
theorem a108211 (n : $\N$) (hn : 0 < n) :\\
\hspace*{2em}(16 * (n : $\Z$) \textasciicircum{} 2 + 1 : $\Z$) =\\
\hspace*{4em}$\lfloor$ 1 / ((4 * (n : $\R$))$^{-1}$ - Real.log 2 +\\
\hspace*{6em}$\sum$ k $\in$ Finset.Icc (n + 1) (2 * n), ((k : $\R$))$^{-1}$) $\rfloor$
\end{quote}
which is, up to the harmless replacement of \texttt{(a n : $\R$) = ($\lfloor \cdot \rfloor$ : $\R$)}
by an equation in $\Z$, character-for-character the statement carried in
\texttt{formal-conjectures} for this conjecture. The sum is over
\texttt{Finset.Icc (n+1) (2*n)}, i.e. $\sum_{k=n+1}^{2n} 1/k$, matching Kimberling's comment.

\emph{A114831.} The \texttt{formal-conjectures} file defines the sequence through
$\Q$-valued division followed by \texttt{Int.toNat} of the floor. We define it instead by
$\N$-valued division,
\begin{quote}\ttfamily\small
a (n+3) = a (n+2) + (2 * a (n+2) * a (n+1)) / (a (n+2) + a (n+1)),
\end{quote}
which is the same function, since truncating division on $\N$ \emph{is} the floor of the rational
quotient of nonnegative integers. The conclusion,
\begin{quote}\ttfamily\small
Tendsto (fun n $\mapsto$ (a (n+1) : $\R$) / (a n : $\R$)) atTop (nhds (Real.sqrt 3)),
\end{quote}
is identical to the corpus statement. The first terms are checked against OEIS by
\texttt{decide}: $a(3)=3$, $a(4)=5$, $a(5)=8$, $a(6)=14$.

\emph{A100434.} Here our statement necessarily differs from the corpus, since the corpus
statement is false; the difference is exactly the minus sign of \eqref{eq:bdef}, and it is
documented in the file header. The development lives in the namespace
\texttt{OeisA100434Corrected} to keep it distinguishable. All other definitions
($a$, $c$, $d$, $e$, $f$, $g$) are transcribed unchanged, including the integer-division form of
$e$ and $f$.

\emph{A114362.} The Lean statement is verbatim \texttt{formal-conjectures}'s \texttt{conjecture2},
including the definition of \texttt{t}:
\begin{quote}\ttfamily\small
theorem conjecture2 :\\
\hspace*{2em}(fun n : $\N$ => (1 - t n) / (1 + t n) -\\
\hspace*{4em}(1 / (2:$\R$)\textasciicircum{}n + 1 / (3:$\R$)\textasciicircum{}n +
1 / (5:$\R$)\textasciicircum{}n + 1 / (7:$\R$)\textasciicircum{}n))\\
\hspace*{2em}=O[atTop] (fun n : $\N$ => 1 / (11:$\R$)\textasciicircum{}n)
\end{quote}
character-for-character identical to the corpus file. Because an \texttt{isBigO} statement only
asserts \emph{some} witness constant, the explicit value is not part of the statement being
formalized; the file supplies $56$ (lemma \texttt{main\_bound}), weaker than the $24$ established
informally in Theorem~\ref{thm:a114362}. The gap is a formalization choice, not a mathematical
one: the Lean proof bounds $|y(n)-S(n)|$ by two applications of the triangle inequality
(\texttt{abs\_sub\_le}, then \texttt{abs\_le}) through the auxiliary quantity
$(1-Q_4(n))/(1+Q_4(n))$, where $Q_4(n)$ is the finite product over $\{2,3,5,7\}$, rather than the
signed one-sided bound $-F \le y(n)-S(n) \le r$ used in the proof above; the two routes give $56$
and $24$ respectively for the same underlying estimates. We did not tighten the Lean constant,
since the corpus statement only asks for \emph{some} $C$.

\subsection{Notable techniques and \texttt{mathlib} workarounds}

\paragraph{The alternating harmonic series is not in \texttt{mathlib}.}
The informal proof of Theorem~\ref{thm:a108211} passes through Mercator's series
$\log 2 = \sum_{j\ge1} (-1)^{j+1}/j$ (Lemma~\ref{lem:mercator}). At the pinned revision,
\texttt{mathlib} has \texttt{Real.hasSum\_log\_sub\_log\_of\_abs\_lt\_one} for $|x| < 1$ but no
statement of the series at the boundary point, and no general alternating-series (Leibniz) test
in \texttt{Analysis/SpecificLimits}. Rather than formalize Abel's theorem or the
integral-remainder argument, the Lean development replaces the whole of Step~1 by an elementary
squeeze, which we consider the most instructive part of the formalization:
\begin{itemize}
\item \texttt{sum\_ffun}: $\sum_{m < n} \frac{1}{(2m+1)(2m+2)} = H_{2n} - H_n$, proved by
induction on $n$. This is Lemmas~\ref{lem:catalan} and~\ref{lem:step1} fused into a
\emph{finite} identity, so the delicate regrouping of a conditionally convergent series never
occurs.
\item \texttt{harm\_upper} and \texttt{harm\_lower}: the two-sided estimate
\[
  \log\frac{2n+1}{n+1} \;\le\; H_{2n} - H_n \;\le\; \log 2 \qquad (n \ge 1),
\]
each obtained by applying \texttt{Real.log\_le\_sub\_one\_of\_pos} to the ratios
$\frac{n+i}{n+i+1}$ and $\frac{n+i+2}{n+i+1}$ respectively and telescoping with
\texttt{Finset.sum\_range\_sub}.
\item Since $\log\frac{2n+1}{n+1} \to \log 2$, the squeeze
(\texttt{tendsto\_of\_tendsto\_of\_tendsto\_of\_le\_of\_le'}) gives
$H_{2n}-H_n \to \log 2$; \texttt{summable\_of\_sum\_range\_le} then gives summability from the
upper bound, and \texttt{hasSum\_iff\_tendsto\_nat\_of\_nonneg} upgrades the partial-sum limit to
\texttt{HasSum ffun (Real.log 2)}.
\end{itemize}
In effect the Lean proof \emph{derives} the alternating harmonic sum as a by-product, using only
$\log x \le x-1$.

\paragraph{Polynomial certificates are discharged by \texttt{linarith}, not \texttt{nlinarith}.}
The five certificates of Appendix~\ref{app:certificates} are degree $5$ to $16$ polynomial
inequalities in one real variable with coefficients up to sixteen digits. General nonlinear
arithmetic does not cope with these. The shift $x = 1+t$, $t \ge 0$ converts each into a claim
that a specific nonnegative combination of the monomials $t^0, \dots, t^{16}$ is positive ---
a \emph{linear} claim in the atoms $t^k$. Each certificate lemma therefore begins
\begin{quote}\ttfamily\small
obtain $\langle$t, ht, rfl$\rangle$ : $\exists$ t, 0 $\le$ t $\wedge$ x = 1 + t :=\\
\hspace*{2em}$\langle$x - 1, by linarith, by ring$\rangle$
\end{quote}
and closes with \texttt{linarith [ht, pow\_nonneg ht 2, \ldots, pow\_nonneg ht 16]}. This is the
single most important implementation decision in \texttt{A108211.lean}: it is what allows a
degree-$16$ certificate with a leading coefficient of $1.1 \times 10^{10}$ to be checked in a few
seconds. The file also raises \texttt{maxHeartbeats} to $2{,}000{,}000$ for the largest of them.

\paragraph{Telescoping to the limit.}
The sandwich of Lemma~\ref{lem:sandwich} is formalized without ever forming the series
$\sum \Delta(m)$: both estimates are proved for every finite partial sum, via the telescoping
lemma \texttt{Finset.sum\_range\_sub'} applied to $m \mapsto h(n+m)$ and $m \mapsto g(n+m)$, and then
transported to the limit with the two \texttt{mathlib} lemmas%
\footnote{\texttt{le\_of\_tendsto\_of\_tendsto'} and \texttt{le\_of\_tendsto'}.}
for comparing limits of sequences with limits of their bounds.
The convergence $h(n+m) \to 0$ needed for the lower estimate comes from the decay certificate
$h(x) \le 1/x$ (Lemma~\ref{lem:decay}, \texttt{cert7} in the file) combined with
\texttt{squeeze\_zero}. Because the window lemmas \texttt{window\_low} and \texttt{window\_high}
are themselves strict, the formal sandwich is stated non-strictly, as
\begin{quote}\ttfamily\small
hf n $\le$ T n $\wedge$ T n $\le$ hf n + gf n,
\end{quote}
and strictness of the final chain is inherited from them; the deficit-retention argument of
Lemma~\ref{lem:sandwich} is therefore present in the paper but not needed in the file.

\paragraph{Floors via Euclidean division.}
In \texttt{A114831.lean} the fractional defect $\theta_n$ of \eqref{eq:ratiorec} is never
introduced. Instead the $\N$-level identity \texttt{Nat.div\_add\_mod} --- $p = m q + r$ with
$0 \le r < m$, for $m = a(k{+}2)+a(k{+}1)$ and $p = 2a(k{+}2)a(k{+}1)$ --- is cast to $\R$ and fed
to a purely algebraic lemma \texttt{ratio\_bound\_core}, which yields the two-sided bound
$F(R_{k+1}) - 1/a(k{+}2) < R_{k+2} \le F(R_{k+1})$ in one step. This is both shorter and more
robust than reasoning about \texttt{Int.floor} of a rational.

\paragraph{A general contraction-with-forcing lemma.}
\texttt{mathlib} has Banach fixed points but nothing directly applicable to a perturbed
contraction $e_{k+1} \le \tfrac12 e_k + b_k$ with $b_k \to 0$ and no a priori bound on
$\sum b_k$. \texttt{A114831.lean} proves the needed statement from scratch
(\texttt{aux\_geom}, \texttt{tendsto\_e\_zero}) by an explicit $\varepsilon$/$N$ argument: given
$b_k \le \varepsilon$ for $k \ge N$, unrolling $m$ steps gives
$e_{N+m} \le 2^{-m} e_N + 2\varepsilon$, and the two terms are then made small separately. This
replaces the tail-splitting computation \eqref{eq:tail} of the paper; the explicit bound
\eqref{eq:a114831bound} is stated in the paper but is not part of the formal statement, which
asserts only the limit.

\paragraph{Integer division and \texttt{omega}.}
In \texttt{A100434.lean} the sequences $e$ and $f$ are $\Z$-valued and defined with integer
division, faithfully to the corpus. The two divisibility lemmas \texttt{d\_even\_even} and
\texttt{d\_odd\_even} ($2 \mid d(2k)$, $2 \mid d(2k+1)$, both by two-step induction) are what make
the halvings exact; once they and the linear identities \eqref{eq:aux5}, \eqref{eq:aux6} are in
context, \texttt{omega} closes both parity cases of \texttt{conjecture2}, since it handles
division by a numeral in the presence of divisibility hypotheses. All inductions use
\texttt{Nat.twoStepInduction}, matching the order-four (parity-decoupled) recurrence.

\paragraph{An elementary Euler-product bridge, avoiding \texttt{tprod} convergence machinery.}
\texttt{A114362.lean} needs $\zeta(m) = \prod_p (1-p^{-m})^{-1}$ for real $m \ge 2$;
\texttt{mathlib} supplies this as \texttt{riemannZeta\_eulerProduct}, a statement about the limit,
along \texttt{Filter.atTop}, of the finite partial products $\mathrm{EP}(m,N) :=
\prod_{p < N,\ p\text{ prime}} (1-p^{-m})^{-1}$ (the primes below $N$ being
\texttt{Nat.primesBelow}~$N$ in the file). Rather than
upgrade this to \texttt{mathlib}'s general infinite-product (\texttt{Multipliable}/\texttt{tprod})
API for products indexed by primes --- machinery not otherwise present in this file, since no
existing lemma there is stated for a product over primes specifically --- the whole development
works with the $N$-indexed finite partial products \texttt{EP} and \texttt{qq} throughout, taking
limits only once, at the very end, in \texttt{tendsto\_prod\_qq}. The peeling of $p=2,3,5,7$ from
the tail is likewise finite: \texttt{Finset.prod\_sdiff} splits \texttt{Nat.primesBelow}~$N$ into
$\{2,3,5,7\}$ and its complement for every fixed $N \ge 8$, and the complement's product is
bounded, for every such $N$, using the finite-product inequality \texttt{one\_sub\_sum\_le\_prod}
and the explicit tail estimate \texttt{tail\_pow} (itself reduced, via a comparison to
$\sum 1/k^2$, to the fully elementary telescoping bound \texttt{tail\_sq}); only the resulting
two-sided bound on the finite product is then passed to the limit, in \texttt{t\_upper} and
\texttt{t\_lower}. This keeps every step of the Euler-product argument at the level of ordinary
\texttt{Finset} sums and products, at the cost of re-deriving, in \texttt{tail\_sq} and
\texttt{tail\_pow}, an elementary series-comparison lemma that a general primes-indexed
\texttt{tprod} API would have supplied directly.

\section{Related work}
\label{sec:related}

Automated and AI-assisted contributions to research mathematics have grown quickly in the last
three years. FunSearch \cite{funsearch} demonstrated that a language model in a genetic
programming loop, with an exact evaluator in the loop, can find record constructions in
extremal combinatorics; AlphaEvolve \cite{alphaevolve} extended the evolutionary-search paradigm
to a broader class of construction and algorithm-design problems. On the community side, the
Erd\H{o}s Problems site and its AI-contributions wiki \cite{erdosproblems} have collected a
substantial number of AI-assisted resolutions and partial results, with an explicit and useful
apparatus of caveats about selection bias, literature coverage, and the difference between an
argument and a verified proof. The \texttt{formal-conjectures} corpus \cite{formalconjectures}
supplies formal statements of open conjectures --- the raw material for work of the kind reported
here --- and \texttt{mathlib} \cite{mathlib} supplies the verified background theory.

Our contribution should be placed modestly relative to that work. The conjectures resolved here
are small; none is a named open problem, and two of the four are elementary. What is
distinctive is the completeness of the chain: statement selection, proof discovery, adversarial
refutation and repair, formalization, and axiom audit were all carried out by automated systems,
and the end product is a machine-checked theorem rather than a natural-language argument
awaiting referee attention. We make no claim about how this scales to harder problems.

The same author who posed A114362 recorded, in the neighbouring entry A348829, both a sharper
even-argument conjecture and a 2024 remark relating the same quantity to the prime zeta function
\cite{ordowski}; Theorem~\ref{thm:a114362}'s Remark ``What is not proved here'' (Section
\ref{sec:a114362}) states precisely how those two further conjectures relate to, and go beyond,
what we prove.

\section{Limitations}

We state the limitations of this paper plainly.

\begin{enumerate}
\item \emph{The conjectures are minor.} All four are conjectural comments in OEIS entries, not
problems of record. A100434 carries OEIS's own \texttt{easy} keyword, and its identities are six
routine inductions. A114831's answer was already guessed correctly in the OEIS entry and in the
\texttt{formal-conjectures} docstring; what was missing, and what we supply, is a proof that the
limit exists at all. Two of the four, Theorems~\ref{thm:a108211} and~\ref{thm:a114362}, required
a genuine idea beyond routine calculation: a rational telescoping certificate accurate to
$O(n^{-7})$ against a $\Theta(n^{-4})$ window in the first case, and an explicit $O(11^{-n})$
error bound extracted from a convergent Euler product in the second; even there the ideas are
classical in kind.

\item \emph{The first attempt was wrong.} Our initial proof of Theorem~\ref{thm:a108211} rested
on a comparison inequality that is false at $m=2$, and would have been insufficiently precise
even if true. The same happened with Theorem~\ref{thm:a114362}: our initial proof asserted
uniform bounds on nested M\"obius-composition corrections that are false already at $n=2$, and
two independent reviewers found the identical defect (Section~\ref{sec:a114362-review}). Both
failures were caught by adversarial review, not by the solver. We report this because the success
rate of the first gate, not of the pipeline, is the honest measure of a solver, and here it is
$2$ out of $4$ theorems, not $4$ out of $4$.

\item \emph{One interpretive step is not mathematics.} The corrected reading \eqref{eq:bdef} of
Dement's $b$ is a judgement about the intent of a natural-language comment. We have set out the
textual and numerical evidence in Section~\ref{sec:a100434}; a reader who rejects it is entitled
to say that we proved a different statement from the one conjectured. The literal statement is
false, so some such judgement is unavoidable.

\item \emph{Literature search was automated.} We checked for prior proofs by automated search and
found none, but automated searches miss things, and OEIS comments of this vintage may well have
been settled in correspondence or in unindexed notes. Priority claims should be read with that
caveat.

\item \emph{Formalization covers the theorems, not the paper.} The Lean files verify the four
theorem statements. They do not verify the explicit error bound \eqref{eq:a114831bound}, the
deficit-retention refinement of Lemma~\ref{lem:sandwich}, the sharper asymptotic of
Theorem~\ref{thm:a114362}'s remark, or the numerical assertions in Section~\ref{sec:upstream};
those were checked by exact rational computation only. The formalized witness constant for
Theorem~\ref{thm:a114362} is $56$, not the $24$ proved informally, for the reason given in
Section~\ref{sec:formal}. Nor did we formalize a corrected version of A109074.

\item \emph{We did not repair the upstream corpus.} The two mis-formalizations of
Section~\ref{sec:upstream} are reported here and left in place; correcting them is a matter
for the maintainers of that corpus.
\end{enumerate}

\appendix
\setcounter{table}{0}
\renewcommand{\thetable}{A.\arabic{table}}

\section{Certificate polynomials for Theorem~\ref{thm:a108211}}
\label{app:certificates}

This appendix makes Section~\ref{sec:a108211} self-contained. Every inequality used there is
reduced to the positivity of a single univariate polynomial with integer coefficients, exhibited
below in the shifted variable $t = x-1 \ge 0$. In each case the polynomial arises as the
numerator of a difference of rational functions over a common denominator that is manifestly
positive for $x \ge 1$ (each denominator is a product of factors of the form $\alpha x + \beta$
with $\alpha, \beta > 0$, and of $16x^2+1$ or $8x^2+1$). Since every listed coefficient is a
positive integer, each numerator is positive for all $t \ge 0$, hence for all real $x \ge 1$; no
approximation, and in particular no decimal, enters at any point.

Recall from \eqref{eq:Apoly}--\eqref{eq:f1f2}:
$A(x) = 1024x^5 + 5888x^4 + 12928x^3 + 13312x^2 + 6148x + 843$,
$B(x) = 4(4x+1)(2x+1)(4x+3)(4x+5)(2x+3)(4x+7)$, $h = A/B$, $g(x) = 60/(4x+1)^7$,
$f_1(x) = \frac{1}{4x} - \frac{1}{16x^2+1}$, $f_2(x) = \frac{1}{4x} - \frac{1}{16x^2+2}$.

\subsection*{C1. The defect identity (Lemma~\ref{lem:defect})}

\[
  \frac{1}{(2m+1)(2m+2)} - \bigl(h(m)-h(m+1)\bigr)
  \;=\; \frac{63\,(16m^2+120m+119)}{2\,P(m)},
\]
\[
  P(m) = (m+1)(2m+1)(2m+3)(2m+5)(4m+1)(4m+3)(4m+5)(4m+7)(4m+9)(4m+11).
\]
This is an identity, not an inequality: after clearing denominators both sides are polynomials of
degree $12$ and the difference is identically zero. It is the only certificate of the five that
is checked by expansion rather than by sign inspection.

\subsection*{C2. The majorant (Lemma~\ref{lem:majorant})}

\[
  g(m) - g(m+1) - \Delta(m) \;=\; \frac{N_2(m)}{D_2(m)},
\]
\[
  D_2(m) = 2\,(m+1)(2m+1)(2m+3)(2m+5)(4m+1)^7(4m+3)(4m+5)^7(4m+7)(4m+9)(4m+11).
\]

\begin{table}[ht]
\centering
\begin{tabular}{@{}rr@{}}
\toprule
$k$ & coefficient of $t^k$ in $N_2(1+t)$ \\
\midrule
14 & 11{,}274{,}289{,}152 \\
13 & 372{,}051{,}542{,}016 \\
12 & 5{,}878{,}132{,}506{,}624 \\
11 & 57{,}656{,}538{,}562{,}560 \\
10 & 387{,}335{,}580{,}549{,}120 \\
 9 & 1{,}873{,}929{,}205{,}972{,}992 \\
 8 & 6{,}717{,}806{,}381{,}039{,}616 \\
 7 & 18{,}122{,}661{,}920{,}047{,}104 \\
 6 & 36{,}993{,}936{,}957{,}112{,}320 \\
 5 & 56{,}922{,}611{,}911{,}495{,}680 \\
 4 & 65{,}063{,}899{,}140{,}917{,}760 \\
 3 & 53{,}630{,}918{,}265{,}563{,}520 \\
 2 & 30{,}166{,}767{,}811{,}516{,}080 \\
 1 & 10{,}374{,}002{,}614{,}521{,}000 \\
 0 & 1{,}646{,}809{,}468{,}266{,}375 \\
\bottomrule
\end{tabular}
\caption{Certificate C2: the shifted numerator $N_2(1+t)$, of degree $14$. Minimum coefficient
$11{,}274{,}289{,}152$; constant term $1{,}646{,}809{,}468{,}266{,}375$. All fifteen coefficients
are positive, so $N_2(1+t) > 0$ for $t \ge 0$.}
\label{tab:C2}
\end{table}

\subsection*{C3. The lower window (Lemma~\ref{lem:window})}

\[
  h(x) - f_1(x) \;=\; \frac{N_3(x)}{D_3(x)},
  \qquad
  D_3(x) = 4x(2x+1)(2x+3)(4x+1)(4x+3)(4x+5)(4x+7)(16x^2+1).
\]

\begin{table}[ht]
\centering
\begin{tabular}{@{}rr@{}}
\toprule
$k$ & coefficient of $t^k$ in $N_3(1+t)$ \\
\midrule
5 & 256 \\
4 & 2{,}816 \\
3 & 12{,}096 \\
2 & 24{,}688 \\
1 & 22{,}223 \\
0 & 6{,}756 \\
\bottomrule
\end{tabular}
\caption{Certificate C3: $N_3(1+t) = 256t^5 + 2816t^4 + 12096t^3 + 24688t^2 + 22223t + 6756$.
Minimum coefficient $256$. Exact value at $x=1$: $h(1)-f_1(1) = 563/294525$.}
\label{tab:C3}
\end{table}

\subsection*{C4. The upper window (Lemma~\ref{lem:window})}

\[
  f_2(x) - h(x) - g(x) \;=\; \frac{N_4(x)}{D_4(x)},
\]
\[
  D_4(x) = 4x(2x+1)(2x+3)(4x+1)^7(4x+3)(4x+5)(4x+7)(8x^2+1).
\]

\begin{table}[ht]
\centering
\begin{tabular}{@{}rr@{}}
\toprule
$k$ & coefficient of $t^k$ in $N_4(1+t)$ \\
\midrule
9 & 393{,}216 \\
8 & 6{,}291{,}456 \\
7 & 43{,}757{,}568 \\
6 & 169{,}863{,}168 \\
5 & 404{,}609{,}280 \\
4 & 614{,}989{,}440 \\
3 & 598{,}254{,}960 \\
2 & 359{,}583{,}120 \\
1 & 120{,}914{,}895 \\
0 & 17{,}203{,}050 \\
\bottomrule
\end{tabular}
\caption{Certificate C4: the shifted numerator $N_4(1+t)$, of degree $9$. Minimum coefficient
$393{,}216$. Exact value at $x=1$: $f_2(1)-h(1)-g(1) = 12743/21656250$.}
\label{tab:C4}
\end{table}

\subsection*{C5. The decay bound (Lemma~\ref{lem:decay})}

\[
  \frac1x - h(x) \;=\; \frac{B(x) - x\,A(x)}{x\,B(x)},
\]
where $x\,B(x) > 0$ for $x \ge 1$ and the numerator, after the shift, is the polynomial of
Table~\ref{tab:C5}.

\begin{table}[ht]
\centering
\begin{tabular}{@{}rr@{}}
\toprule
$k$ & coefficient of $t^k$ in $B(1+t) - (1+t)A(1+t)$ \\
\midrule
6 & 3{,}072 \\
5 & 37{,}120 \\
4 & 184{,}448 \\
3 & 482{,}304 \\
2 & 699{,}916 \\
1 & 534{,}509 \\
0 & 167{,}757 \\
\bottomrule
\end{tabular}
\caption{Certificate C5, establishing $h(x) \le 1/x$ for $x \ge 1$ and hence $h(x) \to 0$.}
\label{tab:C5}
\end{table}

\medskip\noindent
All five certificates were computed by exact rational arithmetic, recomputed independently by a
reviewing model, and re-derived a third time inside Lean, where C2--C5 appear as the lemmas
\texttt{cert2}, \texttt{cert5}, \texttt{cert6}, \texttt{cert7} and C1 as \texttt{delta\_eq}.

\begin{thebibliography}{99}

\bibitem{oeis}
OEIS Foundation Inc., \emph{The On-Line Encyclopedia of Integer Sequences}.
\url{https://oeis.org}. Sequences cited: A100434 (N.~J.~A.~Sloane, 21 November 2004, suggested by
correspondence from C.~Dement; conjectural comment of C.~Dement, 18 December 2004),
A108211 (R.~Zumkeller, 15 June 2005; conjectural comment of C.~Kimberling, 9 September 2014),
A114831 (J.~V.~Post, 19 February 2006), A114362 (B.~Cloitre, 9 February 2006; conjectural
comments of T.~Ordowski, 5 January 2022 and 13 November 2022), A348829 (T.~Ordowski,
1 November 2021; conjectural comments of 13 November 2022 and 6 November 2024). Also
referenced: A001333, A001764, A002315, A005156, A005319, A052542, A075870.

\bibitem{kimberling}
C.~Kimberling, comment on OEIS A108211, 9 September 2014:
``Conjecture: $a(n) = \lfloor 1/(1/(4n) - \log(2) + 1/(n+1) + 1/(n+2) + \cdots + 1/(2n)) \rfloor$.''

\bibitem{dement}
C.~Dement, comment on OEIS A100434, 18 December 2004, defining the auxiliary sequences
$b, c, d, e, f, g$ and conjecturing
$c(n)+d(n) = e(n)+f(n) = g(n)+a(n) = b(n)$.

\bibitem{ordowski}
T.~Ordowski, comments on OEIS A114362 and A348829. A114362, 13 November 2022:
``Conjecture: $(1-t(n))/(1+t(n)) = 1/2^n+1/3^n+1/5^n+1/7^n+O(1/11^n)$, where
$t(n)=\zeta(2n)/\zeta(n)^2$.'' A348829, 13 November 2022: ``Conjecture:
$0 < w(2n) - (1/2^{2n}+1/3^{2n}+1/5^{2n}+1/7^{2n}) < 1/11^{2n}$ for every $n>0$.'' A348829,
6 November 2024: ``It can be proven that $P(2n) - w(2n) \sim 1/12^{2n}$, where
$P(x) = \sum_{p} 1/p^x$ is the prime zeta function of real $x>1$.''

\bibitem{formalconjectures}
Google DeepMind, \emph{formal-conjectures}: a collection of formalized conjectures in Lean~4.
\url{https://github.com/google-deepmind/formal-conjectures}. Files referenced:
\texttt{FormalConjectures/OEIS/100434.lean}, \texttt{108211.lean}, \texttt{109074.lean},
\texttt{114362.lean}, \texttt{114831.lean}.

\bibitem{lean4}
L.~de Moura and S.~Ullrich, \emph{The Lean 4 theorem prover and programming language},
in: Automated Deduction --- CADE 28, Lecture Notes in Computer Science 12699, Springer, 2021,
pp.~625--635.

\bibitem{mathlib}
The mathlib Community, \emph{The Lean mathematical library},
in: Proceedings of the 9th ACM SIGPLAN International Conference on Certified Programs and Proofs
(CPP 2020), ACM, 2020, pp.~367--381.
\url{https://github.com/leanprover-community/mathlib4}.

\bibitem{funsearch}
B.~Romera-Paredes, M.~Barekatain, A.~Novikov, M.~Balog, M.~P.~Kumar, E.~Dupont, F.~J.~R.~Ruiz,
J.~S.~Ellenberg, P.~Wang, O.~Fawzi, P.~Kohli, and A.~Fawzi,
\emph{Mathematical discoveries from program search with large language models},
Nature \textbf{625} (2024), 468--475.

\bibitem{alphaevolve}
A.~Novikov et al., \emph{AlphaEvolve: a coding agent for scientific and algorithmic discovery},
Google DeepMind, 2025.

\bibitem{erdosproblems}
T.~F.~Bloom (site), T.~Tao (community database), and contributors, \emph{Erd\H{o}s Problems} and the associated
AI-contributions wiki. \url{https://www.erdosproblems.com},
\url{https://github.com/teorth/erdosproblems}.

\end{thebibliography}

\end{document}
